%!TEX root = ../main.tex
\section{Заключение}

\subsection{Идеи на будущее}

Если бы у меня было ещё время на развитие проекта, я бы с удовольствием реализовал следующие идеи:
\begin{enumerate}[leftmargin=*,noitemsep]
    \item \textbf{Ускорение построения графа}: с помощью диаграмм алгоритмов кластеризации можно проредить 2000 ориентиров до 300 ключевых развилок без потери точности, ускорив построение графа в 40 раз.
    \item \textbf{Ансамблирование окон для AntMaze Giant}: комбинировать узкое окно ($W=4$) для щелей и широкое ($W=16$) для открытых залов.
    \item \textbf{дополнительное исследование проблем}: Все алгоритмы до сих пор хоть и строят отличный маршрут от старта до конца при рельном прогоне имеют петли и топтание на месте. Нужно дополнительно исследовать причины и то как с этим бороться.
\end{enumerate}

\subsection{Главные выводы}

В ходе этой работы мы прошли полный исследовательский цикл:
\begin{itemize}[leftmargin=*,noitemsep]
    \item Убедились, почему простое деление пополам (Recursive Bisection) ломается в невыпуклых тупиках.
    \item Построили надёжный топологический граф Search on the Replay Buffer со скользящим окном упреждения ($L=2.6\,\text{м}$).
    \item Успешно дистиллировали Дейкстру в нейросеть, ускорив прямой проход до \textbf{0.55\,мс} при точности \textbf{86.0\%} на Medium.
    \item Создали метод \textbf{Enhanced Sequence Waypoint Attention Transformer}, который за счёт токенов кривизны, штрафа ALiBi и окна $4+1$ существенно превзошёл бейзлайн на лабиринте Large (\textbf{60.9\% против 49.3\%}).
\end{itemize}

Главный вывод: Forward-Backward представления содержат богатейшую информацию о геометрии среды, а правильное рассуждение о последовательностях интенций позволяет решать сложнейшие задачи zero-shot навигации легко, элегантно и без тяжёлых генеративных моделей мира.
